\renewcommand{\thetheorem}{14.\arabic{theorem}}
\setcounter{theorem}{0}
\setcounter{chapter}{13} % Sets chapter counter so the current chapter is 14

% ---------------------------------------------------------------------------
% Provenance of the prose in this file.
%   % Extractor: <model>   the block below was transcribed from Prof. Biran's
%                          handwritten notes by that model
%   % Creator:   <model>   the block below was written by that model and has no
%                          counterpart in the notes
% A tag holds until the next tag.
% ---------------------------------------------------------------------------

\chapter{Sums of Vector Spaces}

% Extractor: Gemini 3.1 Pro
% def:14.1
\begin{definition}[Sum of Subspaces]
\label{def:sum_of_subspaces}
Let $V$ be a vector space over $K$, and let $U, W \subseteq V$ be linear subspaces. We define the \newterm{sum} of $U$ and $W$, denoted by $U + W$, as the set:
\[
U + W := \{u + w \mid u \in U, w \in W\}.
\]
In a similar way, we can define sums of more than two subspaces. If $U_1, \dots, U_r \subseteq V$ are subspaces, we define their total sum as:
\[
U_1 + \dots + U_r := \{u_1 + \dots + u_r \mid u_1 \in U_1, \dots, u_r \in U_r\}.
\]
\end{definition}

% prop:14.2
\begin{proposition}[Properties of Sums]
\label{prop:properties_of_sums}
Let $U, W \subseteq V$ be subspaces. Then, the following properties hold:
\begin{enumerate}[label=\textbf{(\alph*)}]
    \item $U + W = \Sp(U \cup W)$. In particular, $U + W \subseteq V$ is a subspace, and $U$ and $W$ are themselves subspaces of $U + W$.
    \item If $U$ and $W$ are finite-dimensional, then so is $U + W$. Moreover, one can systematically find a basis for $U + W$ by following these sequential steps: first, choose a basis $(v_1, \dots, v_k)$ for $U \cap W$ (where $k = \dim(U \cap W)$). Second, extend this to a basis $(v_1, \dots, v_k, u_1, \dots, u_{\ell-k})$ for $U$ (where $\ell = \dim U$). Third, extend the initial intersection basis to a basis $(v_1, \dots, v_k, w_1, \dots, w_{m-k})$ for $W$ (where $m = \dim W$). Finally, the combined elements 
    \[(v_1, \dots, v_k, u_1, \dots, u_{\ell-k}, w_1, \dots, w_{m-k})\] 
    form a basis for $U + W$.
    \item \label{item:dimension_formula_sums} If $U$ and $W$ are finite-di\-men\-sional, the dimensions satisfy the \newterm{dimension formula}
    \[
    \dim(U + W) = \dim U + \dim W - \dim(U \cap W).
    \]
\end{enumerate}
\end{proposition}

\begin{proof}
% Extractor: Opus 5
\textbf{(a)} Clearly $U + W \subseteq \Sp(U \cup W)$, by the characterisation of $\Sp(U \cup W)$ as the set of all linear combinations of elements from $U \cup W$ (\cref{lem:span_is_linear_combinations}, applicable since $U \cup W$ contains $0_V$ and is in particular non-empty). But we also have $\Sp(U \cup W) \subseteq U + W$. Indeed, if $v \in \Sp(U \cup W)$, then
\[
v = \underbrace{\sum_{i=1}^{s} a_i u_i}_{\in U} + \underbrace{\sum_{j=1}^{r} b_j w_j}_{\in W},
\]
with $a_i, b_j \in K$, $u_i \in U$ and $w_j \in W$, which implies that $v \in U + W$. This proves $\Sp(U \cup W) \subseteq U + W$. Putting all together, we have $U + W = \Sp(U \cup W)$. The rest of the statements in \textbf{(a)} follow immediately.

% Extractor: Gemini 3.1 Pro
\textbf{(b)} That the subspace $U + W$ is generated by the combined vectors
\[
(v_1, \dots, v_k, u_1, \dots, u_{\ell-k}, w_1, \dots, w_{m-k})
\]
is not quite immediate from the definitions, and the notes leave it to the reader, so it is set as \cref{exc:combined_family_spans} below and solved at the end of the chapter. It remains to show that these vectors are linearly independent.

Suppose there exists a linear combination equal to the zero vector:
\begin{equation}
\label{eq:linear_comb_sum}
\alpha_1 v_1 + \dots + \alpha_k v_k + \beta_1 u_1 + \dots + \beta_{\ell-k} u_{\ell-k} + \gamma_1 w_1 + \dots + \gamma_{m-k} w_{m-k} = 0.
\end{equation}
We can rearrange this equation to isolate the terms belonging to $U$ and $W$:
\[
\underbrace{\beta_1 u_1 + \dots + \beta_{\ell-k} u_{\ell-k}}_{\in U} = \underbrace{-\alpha_1 v_1 - \dots - \alpha_k v_k - \gamma_1 w_1 - \dots - \gamma_{m-k} w_{m-k}}_{\in W}.
\]
Let us denote this common vector by $x$. Since the left side is entirely in $U$ and the right side is entirely in $W$, we have $x \in U$ and $x \in W$, which implies $x \in U \cap W$. 

Because $(v_1, \dots, v_k)$ is a basis for $U \cap W$, there must exist scalars $\delta_1, \dots, \delta_k \in K$ such that:
\[
x = \delta_1 v_1 + \dots + \delta_k v_k.
\]
Equating this with our original expression for $x$ in $U$, we obtain:
\[
\beta_1 u_1 + \dots + \beta_{\ell-k} u_{\ell-k} = \delta_1 v_1 + \dots + \delta_k v_k,
\]
or, after moving everything to one side,
\[
\sum_{i=1}^{\ell-k} \beta_i u_i - \sum_{j=1}^k \delta_j v_j = 0.
\]
However, the vectors $(v_1, \dots, v_k, u_1, \dots, u_{\ell-k})$ form a basis for $U$, meaning they are strictly linearly independent. Consequently, all coefficients in this sum must be zero. In particular, $\beta_1 = \dots = \beta_{\ell-k} = 0$.

Substituting these zero values back into our initial equation \eqref{eq:linear_comb_sum}, the terms involving $u_i$ vanish, leaving:
\[
\alpha_1 v_1 + \dots + \alpha_k v_k + \gamma_1 w_1 + \dots + \gamma_{m-k} w_{m-k} = 0.
\]
Since the vectors $(v_1, \dots, v_k, w_1, \dots, w_{m-k})$ form a basis for $W$, they are linearly independent. This immediately implies that $\alpha_1 = \dots = \alpha_k = 0$ and $\gamma_1 = \dots = \gamma_{m-k} = 0$. 

Because all coefficients $\alpha_i$, $\beta_i$, and $\gamma_j$ have been forced to zero, the combined family of vectors is linearly independent, and therefore constitutes a valid basis for $U + W$.

Statement \textbf{(c)} now follows by simply counting the basis vectors produced in \textbf{(b)}. The basis consists of the $k$ vectors $v_i$, the $\ell - k$ vectors $u_i$, and the $m - k$ vectors $w_j$, and therefore,
\begin{align*}
\dim(U + W) &= k + (\ell - k) + (m - k) \\
&= \ell + m - k \\
&= \dim U + \dim W - \dim(U \cap W).
\end{align*}
\end{proof}

% Extractor: Opus 5
\begin{exercise}[The Combined Family Spans the Sum]
\label{exc:combined_family_spans}
The proof above verifies that the family
\[
(v_1, \dots, v_k, u_1, \dots, u_{\ell-k}, w_1, \dots, w_{m-k})
\]
is linearly independent. Check that it also spans $U + W$, which the notes leave to the reader.
\end{exercise}

% Extractor: Gemini 3.1 Pro
% cor:14.3
\begin{corollary}[Equivalent Conditions for Trivial Intersection]
\label{cor:trivial_intersection}
% Extractor: Opus 5
Let $V$ be a vector space, and let $U, W \subseteq V$ be two finite-di\-men\-sional linear subspaces. The following statements are equivalent:
\begin{enumerate}[label=\textbf{(\alph*)}]
    \item $\dim(U + W) = \dim U + \dim W$.
    \item $\dim(U \cap W) = 0$.
% Extractor: Gemini 3.1 Pro
    \item $U \cap W = \{0_V\}$.
    \item For every $v \in U + W$, there exist unique vectors $u \in U$ and $w \in W$ such that $v = u + w$.
    \item The equality $0 = u + w$ with $u \in U$ and $w \in W$ implies that $u = 0$ and $w = 0$.
\end{enumerate}
\end{corollary}

\begin{proof}
% Extractor: Opus 5
The equivalence \textbf{(a)} $\iff$ \textbf{(b)} follows from the previous \cref{prop:properties_of_sums} \textbf{(c)}, whose dimension formula reads $\dim(U + W) = \dim U + \dim W - \dim(U \cap W)$. And \textbf{(b)} $\iff$ \textbf{(c)} holds because there is only one linear subspace of dimension $0$, namely the space $\{0_V\}$: a subspace of dimension $0$ has the empty set as a basis, and $\Sp(\emptyset) = \{0_V\}$ by \cref{rem:span_of_empty_set}, which is the reading of \cref{rem:dimension_of_zero_space} in the other direction.

The remaining three statements are tied to \textbf{(c)} by the cycle of implications \textbf{(c)} $\implies$ \textbf{(d)} $\implies$ \textbf{(e)} $\implies$ \textbf{(c)}.

% Extractor: Gemini 3.1 Pro
\textbf{(c) $\implies$ (d):} Suppose that $v \in U + W$ can be written in two ways: $v = u + w = u' + w'$, where $u, u' \in U$ and $w, w' \in W$. Rearranging this equation yields $u - u' = w' - w$. The left side is in $U$ and the right side is in $W$, which means the vector $u - u'$ resides in the intersection $U \cap W$. Since we assume $U \cap W = \{0_V\}$, it follows that $u - u' = 0$ and $w' - w = 0$. Therefore, $u = u'$ and $w = w'$, proving uniqueness.

\textbf{(d) $\implies$ (e):} Consider the zero vector, which can obviously be written as $0 = 0_U + 0_W$. If we also have $0 = u + w$ for some $u \in U$ and $w \in W$, the assumed uniqueness in statement \textbf{(d)} immediately forces $u = 0$ and $w = 0$.

\textbf{(e) $\implies$ (c):} Let $x \in U \cap W$. We can express the zero vector as $0 = x + (-x)$. Since $x \in U \cap W$, we know $x \in U$ and $-x \in W$. If statement \textbf{(e)} holds, this equality strictly implies $x = 0$. Consequently, the only element in the intersection is the zero vector, so $U \cap W = \{0_V\}$.
\end{proof}

% def:14.4
\begin{definition}[Complement of a Subspace]
\label{def:complement}
Let $V$ be a vector space over $K$, and let $U \subseteq V$ be a subspace. A subspace $W \subseteq V$ is called a \newterm{complement} of $U$ if it satisfies the following two conditions:
\[
U + W = V \quad \text{and} \quad U \cap W = \{0_V\}.
\]
\end{definition}

\begin{remark}
Note that if $W$ is a complement of $U$, then $U$ and $W$ satisfy all of the five equivalent statements derived from \cref{cor:trivial_intersection}.
\end{remark}

% prop:14.5
\begin{proposition}[Existence of a Complement]
\label{prop:existence_complement}
Let $V$ be a finite-dimensional vector space, and let $U \subseteq V$ be a subspace. Then, there exists a subspace $W \subseteq V$ which is a complement to $U$.
\end{proposition}

\begin{proof}
Note first that $U$ is finite-di\-men\-sional, by \cref{prop:subspace_dimension_constraints}. Choose a basis $\mathcal{B}_U = (u_1, \dots, u_{\ell})$ for $U$, where $\ell = \dim U$. By the basis extension property, item \textbf{(2)} of \cref{sum:finite_dim_properties}, we can extend this to a full basis for $V$:
\[
\mathcal{B}_V = (u_1, \dots, u_{\ell}, w_1, \dots, w_m),
\]
where $\ell + m = \dim V$. We define $W := \Sp(w_1, \dots, w_m)$. 

We claim that $W$ is a valid complement of $U$. Indeed, since the combined basis spans the entirety of $V$, it is clear that $U + W = V$: by \cref{prop:properties_of_sums} \textbf{(a)}, $U + W = \Sp(U \cup W)$, and this span contains all of $u_1, \dots, u_{\ell}, w_1, \dots, w_m$, hence contains $\Sp(\mathcal{B}_V) = V$. To see that $U \cap W = \{0\}$, let $x \in U \cap W$. Because $x \in U$, we can write $x = \sum_{i=1}^{\ell} a_i u_i$. Because $x \in W$, we can also write $x = \sum_{j=1}^m b_j w_j$. Equating these expressions yields:
\[
\sum_{i=1}^{\ell} a_i u_i = \sum_{j=1}^m b_j w_j \implies \sum_{i=1}^{\ell} a_i u_i - \sum_{j=1}^m b_j w_j = 0.
\]
Since $\mathcal{B}_V$ is a basis for $V$, all of its constituent vectors are strictly linearly independent. This forces all scalar coefficients to be zero: $a_i = 0$ for all $i$ and $b_j = 0$ for all $j$. Thus, $x = 0$, proving that $U \cap W = \{0\}$.
\end{proof}

% Extractor: Opus 5
\begin{importantremark}
\label{rem:complement_not_unique}
If $\{0_V\} \subsetneq U \subsetneq V$, then there are many different subspaces $W \subseteq V$ which are complements of $U$. In general, there is no canonical choice for a complement $W$. (Both strict inclusions are needed. The subspace $U = \{0_V\}$ has the single complement $W = V$, since $\{0_V\} + W = W$, and $U = V$ has the single complement $W = \{0_V\}$; in either case there is nothing to choose.)
\end{importantremark}

% Creator: Opus 5, TikZ rendering of the sketch in the notes
\begin{figure}[htbp]
  \centering
  \begin{tikzpicture}[scale=1.0, font=\small]
    % The ambient plane, drawn only as a faint pair of coordinate axes: the
    % three subspaces below are the actual content of the picture.
    \draw[gray!30, thin] (-2.9,0) -- (2.9,0);
    \draw[gray!30, thin] (0,-2.3) -- (0,2.3);

    %, U, the subspace to be complemented ---
    \draw[ThemeSecondary, very thick] (-2.8,-1.0) -- (2.8,1.0);
    \node[ThemeSecondary, anchor=west, inner sep=3pt] at (2.8,1.0) {$U$};

    %, two different complements of U ---
    \draw[ThemeGreen, very thick] (-1.35,-2.2) -- (1.35,2.2);
    \node[ThemeGreen, anchor=south, inner sep=3pt] at (1.35,2.2) {$W$};

    \draw[ThemePurple, very thick] (-2.4,2.1) -- (2.4,-2.1);
    \node[ThemePurple, anchor=south east, inner sep=3pt] at (-2.4,2.1) {$W'$};

    \fill[black] (0,0) circle (1.7pt);
    \node[anchor=north west, inner sep=4pt] at (0.15,-0.1) {$0$};
    \node[gray!60!black, anchor=north east] at (2.9,-0.15) {$V = \mathbb{R}^{2}$};
  \end{tikzpicture}
  \caption{Two complements of the same subspace. Inside $V = \mathbb{R}^{2}$,
  take for $U$ any line through the origin. Every other line through the origin
  meets $U$ only in $0$ and together with $U$ spans the whole plane, so both $W$
  and $W'$ are complements of $U$, and so is any of the infinitely many
  further lines through the origin.}
  \label{fig:complements_not_unique}
\end{figure}

% Creator: Opus 5
\subsection{Solutions to the Exercises}

\begin{exercisesolution}[Proof that the Combined Family Spans $U + W$ (on \cpageref{exc:combined_family_spans})]
Write $\mathcal{S} := (v_1, \dots, v_k, u_1, \dots, u_{\ell-k}, w_1, \dots, w_{m-k})$ for the combined family.

Every member of $\mathcal{S}$ lies in $U + W$: the vectors $v_i$ and $u_i$ lie in $U$, the vectors $w_j$ lie in $W$, and both $U$ and $W$ are contained in $U + W$ by \cref{prop:properties_of_sums} \textbf{(a)}. Since $U + W$ is a subspace, it follows that $\Sp(\mathcal{S}) \subseteq U + W$.

For the reverse inclusion, let $x \in U + W$ and write $x = u + w$ with $u \in U$ and $w \in W$. Because $(v_1, \dots, v_k, u_1, \dots, u_{\ell-k})$ is a basis for $U$, there are scalars with
\[
u = \sum_{i=1}^{k} \alpha_i v_i + \sum_{i=1}^{\ell-k} \beta_i u_i,
\]
and because $(v_1, \dots, v_k, w_1, \dots, w_{m-k})$ is a basis for $W$, there are scalars with
\[
w = \sum_{i=1}^{k} \alpha'_i v_i + \sum_{j=1}^{m-k} \gamma_j w_j.
\]
Adding the two expressions writes $x$ as a linear combination of the members of $\mathcal{S}$, the coefficient of $v_i$ being $\alpha_i + \alpha'_i$. Hence $U + W \subseteq \Sp(\mathcal{S})$, and the two inclusions give $\Sp(\mathcal{S}) = U + W$.
\end{exercisesolution}

\begin{ainote}
\Cref{cor:trivial_intersection} deserves a word. The notes state \emph{five}
equivalent conditions and the transcript kept only four: the two dimension
criteria, $\dim(U + W) = \dim U + \dim W$ and $\dim(U \cap W) = 0$, were lost,
and in their place stood a garbled duplicate of statement \textbf{(e)}. Prof.
Biran's list is restored, together with his hypothesis that $U$ and $W$ be
finite-di\-men\-sional, without which the two dimension criteria mean nothing.
That something had gone missing was visible in the transcript itself, whose
remark after \cref{def:complement} already spoke of \qt{the five equivalent
statements}.

One correction of statement. \Cref{rem:complement_not_unique} said that a
subspace $U \subsetneq V$ has many complements. It needs $U \ne \{0_V\}$ as well:
the zero subspace has $W = V$ as its only complement, since $\{0_V\} + W = W$,
and there is nothing to choose. The figure beside the remark had assumed as much
all along, taking for $U$ a line rather than a point, which is the sort of
disagreement between a picture and its caption worth checking for elsewhere.
\end{ainote}
